\begin{table}[!htbp]\centering
\caption{Calibrated price-discovery effect vs.\ the data's ceiling}
\label{tab:calibration}
\resizebox{\textwidth}{!}{%
\begin{tabular}{lccc}
\toprule
 & \multicolumn{3}{c}{Buyer-error share of residual variance}\\
\cmidrule(lr){2-4}
Comp window & 0.10 & 0.33 & 1.00\\
\midrule
$k=4$ & 0.009 & 0.055 & 0.226\\
$k=6$ & 0.013 & 0.082 & 0.338\\
$k=8$ & 0.017 & 0.109 & 0.449\\
\midrule
\multicolumn{4}{l}{Data: 95\% one-sided ceiling (IQR measure) = 0.17 log points}\\
\multicolumn{4}{l}{Pinned $\sigma_a$ (Zillow within-ZIP-year residual SD) = 0.188; published off-market error $\approx$ 0.10}\\
\bottomrule
\end{tabular}}
\begin{minipage}{0.92\linewidth}\footnotesize\vspace{2pt}
\emph{Notes:} Each cell is the model's predicted rise in within-segment
residual dispersion (log points, at mean Zillow exposure) when the machine
signal is removed, $\rho(v_m)(V_0-V_1)$ converted to a standard-deviation
change, evaluated at the machine's deed footprint ($v_m \approx$ twice the
mean purchase share) and $\sigma_a$ pinned from Zillow's own purchase-price
residual dispersion. The data's 95\% one-sided ceiling (
0.17 log points) lies \emph{below} the model's
most aggressive cells (large comp windows and high buyer-error share; 3 of 9 columns shown), which the design therefore rejects, and above its
central calibrations, with which it is consistent. The externality is
bounded tightly, not point-identified; either way it is economically small.
\end{minipage}
\end{table}
